% Apunts de teoria en PDF · 3r ESO · Unitat 10: Proporcionalitat. Teorema de Tales
% Mateix estil que els apunts de Batxillerat (franja i caixes de teoria en blau, exemples en roig), però més visuals,
% curts i amb dibuixos, per a l'alumnat de 3r ESO. Es compila amb pdflatex dues vegades (per als enllaços de l'índex),
% amb logo-fentmates.pdf (_eines/logo/pdf/horitzontal-clar.pdf) al costat.
\documentclass[11pt,a4paper]{article}
\usepackage[utf8]{inputenc}
\usepackage[T1]{fontenc}
\usepackage[catalan]{babel}
\usepackage[sfdefault,book,lining]{FiraSans}
\usepackage[italic,noplusnominus,noequal]{mathastext}
\linespread{1.18}
\usepackage[a4paper,top=2.4cm,bottom=2.2cm,left=1.6cm,right=1.6cm,headheight=1.1cm,headsep=0.5cm,footskip=1cm]{geometry}
\usepackage{amsmath,amssymb,cancel}
\usepackage{icomma}
\usepackage{graphicx}
\usepackage[table]{xcolor}
\usepackage{tikz}\usetikzlibrary{calc,positioning,arrows.meta,decorations.pathreplacing,shapes.geometric}
\usepackage{fontawesome5}
\usepackage{ccicons}
\usepackage{fancyhdr}
\usepackage{enumitem}
\usepackage{array}
\usepackage{needspace}
\usepackage[most]{tcolorbox}
\usepackage{titlesec}
\usepackage[hidelinks]{hyperref}
\hypersetup{pdftitle={Apunts de teoria. Unitat 10: Proporcionalitat. Teorema de Tales. 3r ESO},pdfsubject={Matemàtiques 3r ESO},pdfkeywords={raó, proporció, teorema de Tales, posició de Tales, semblança, raó de semblança, criteris de semblança, ombres, escales}}

% colors de la web (bloc d'Àlgebra: blau)
\definecolor{fmnavy}{HTML}{1B2748}
\definecolor{fmblue}{HTML}{3D6AA0}
\definecolor{fmbluel}{HTML}{4E81BD}
\definecolor{fmyellow}{HTML}{FFC94D}
\definecolor{fmgold}{HTML}{E9A919}
\definecolor{fmaccent}{HTML}{FF7A50}
\definecolor{fmdef}{HTML}{1D4F73}\definecolor{fmdefbg}{HTML}{EEF4F9}
\definecolor{fmex}{HTML}{A3203A}\definecolor{fmexbg}{HTML}{FBF0F2}
\definecolor{fmsf}{HTML}{C2410C}
\definecolor{fmok}{HTML}{15803D}
\definecolor{fmmuted}{HTML}{5D6685}
\definecolor{fmsurface}{HTML}{EEF1F8}
\definecolor{fmA}{HTML}{2F6FD6}\definecolor{fmB}{HTML}{EA6A10}\definecolor{fmR}{HTML}{E11D48}\definecolor{fmK}{HTML}{16A34A}\definecolor{fmV}{HTML}{7C3AED}

% lletres
\newcommand{\comic}{\fontfamily{ComicNeue-TLF}\selectfont}
\newcommand{\heavy}{\fontseries{eb}\selectfont}
\newcommand{\semi}{\fontseries{sb}\selectfont}
\DeclareUnicodeCharacter{2212}{-}

\setlength{\parindent}{0pt}
\setlength{\parskip}{7pt}
\setlength{\lineskip}{4pt}\setlength{\lineskiplimit}{4pt}
\setlength{\emergencystretch}{3em}
\setlist{itemsep=3pt,topsep=3pt,leftmargin=1.4em}
\setcounter{secnumdepth}{1}
\renewcommand{\thesection}{10.\arabic{section}}
\titleformat{\section}[hang]{\color{fmnavy}\fontsize{23}{27}\heavy\selectfont}{\tcbox[on line,colback=fmblue,colframe=fmblue,arc=3mm,boxsep=2pt,left=6pt,right=6pt,top=3pt,bottom=3pt]{\color{white}\thesection}}{0.45em}{}[\vspace{5pt}{\color{fmblue}\titlerule[2pt]}]
\titlespacing*{\section}{0pt}{4pt}{6pt}
\titleformat{\subsection}{\color{fmblue}\Large\semi}{}{0pt}{\faCaretRight\ }
\titlespacing*{\subsection}{0pt}{8pt}{2pt}

% frase de cada apartat, «a mà»
\newcommand{\frase}[1]{\par\vspace{-2pt}{\comic\Large\color{fmbluel}\faCommentDots[regular]\ \ #1\par}\vspace{2pt}}
\newcommand{\seccio}[3]{\Needspace{16\baselineskip}\section{#2}\hypertarget{#1}{}\frase{#3}}
% notes a mà dins dels exemples
\newcommand{\nota}[1]{{\comic\small\color{fmex!85!black}#1}}
\newcommand{\notab}[1]{{\comic\small\color{fmblue}#1}}
% bé / malament
\newcommand{\be}[1]{{\color{fmok}\faCheckCircle}\ \ensuremath{#1}}
\newcommand{\mal}[1]{{\color{fmR}\faTimesCircle}\ {\color{fmR}\ensuremath{#1}}}
% pas numerat
\newcommand{\pas}[1]{\tikz[baseline=(p.base)]\node[circle,fill=fmblue,text=white,inner sep=1.2pt,minimum size=1.25em,font=\small\heavy](p){#1};}
% etiqueta de color
\newcommand{\chip}[2]{\tcbox[on line,colback=#1!14,colframe=#1!60,arc=2mm,boxsep=1pt,left=5pt,right=5pt,top=1pt,bottom=1pt]{#2}}
\newcommand{\fmkey}[1]{\colorbox{fmyellow!55}{\strut #1}}
\newcommand{\cA}[1]{{\color{fmA}#1}}\newcommand{\cB}[1]{{\color{fmB}#1}}\newcommand{\cK}[1]{{\color{fmK}#1}}\newcommand{\cR}[1]{{\color{fmR}#1}}\newcommand{\cV}[1]{{\color{fmV}#1}}

% caixes
\tcbset{fmbase/.style={enhanced,arc=3mm,boxrule=1.2pt,left=4mm,right=4mm,top=2mm,bottom=2.5mm,before skip=8pt,after skip=8pt,fonttitle=\heavy\large,before upper=\raggedright}}
\newtcolorbox{idea}[1][Idea clau]{fmbase,colback=fmdefbg,colframe=fmdef,coltitle=white,title={\faKey\ \ #1}}
\newtcolorbox{exemple}[1][Exemple]{fmbase,colback=fmexbg,colframe=fmex,colbacktitle=fmexbg,coltitle=fmex,title={\faPencil*\ \ #1},titlerule=0.5pt,titlerule style={fmex!40,dashed}}
\newtcolorbox{compte}[1][Compte!]{fmbase,colback=fmsf!6,colframe=fmsf,coltitle=white,fonttitle=\comic\bfseries\Large,title={\faExclamationTriangle\ \ #1}}
\newtcolorbox{truc}[1][Truc]{fmbase,colback=fmyellow!20,colframe=fmgold,colbacktitle=fmyellow,coltitle=fmnavy,fonttitle=\comic\bfseries\Large,title={\faLightbulb[regular]\ \ #1}}

% dibuix de la targeta de la unitat, en blanc
\newcommand{\fmcard}{\begin{tikzpicture}[x=0.062cm,y=-0.062cm,white,line width=3pt,line cap=round,line join=round]
  \begin{scope}[transparency group,opacity=0.15]
  \node[anchor=base west,inner sep=0pt,font=\fontsize{32}{32}\selectfont] at (14,44) {$\frac{a}{b}=\frac{a'}{b'}$};
  \draw (130,110) -- (225,25); \draw (130,110) -- (232,80);
  \draw[line width=2pt] (165,60) -- (176,96); \draw[line width=2pt] (200,40) -- (215,85);
  \end{scope}
\end{tikzpicture}}
\newcommand{\fmbrand}{\raisebox{-0.25\height}{\includegraphics[height=0.5cm]{logo-fentmates.pdf}}}

% regla de Ruffini: #1 = a, #2 = coeficients, #3 = fila del mig (sense el primer), #4 = fila de baix, #5 = nombre de coeficients
\newcommand{\ruffini}[5]{\begin{tikzpicture}[x=1.25cm,y=0.85cm,every node/.style={font=\large}]
  \foreach \v [count=\i] in {#2} \node (t\i) at (\i,0) {$\v$};
  \foreach \v [count=\i] in {#3} {\pgfmathtruncatemacro\j{\i+1}\node (m\j) at (\j,-1) {$\v$};}
  \foreach \v [count=\i] in {#4} \node (b\i) at (\i,-2.1) {$\v$};
  \node[text=fmB] at (0,-1) {$#1$};
  \draw[line width=1.1pt,fmnavy] (0.45,0.5)--(0.45,-1.5) (0.45,-1.5)--(#5+0.5,-1.5);
  \node[draw=fmR,line width=1.3pt,rounded corners=2mm,minimum size=0.75cm] at (#5,-2.1) {};
  \draw[-{Latex[length=2.2mm]},fmK,line width=1.1pt] (t1.south) -- (b1.north);
  \pgfmathtruncatemacro\last{#5-1}
  \foreach \i in {1,...,\last} {\pgfmathtruncatemacro\j{\i+1}
    \draw[-{Latex[length=2.2mm]},fmB,line width=1.1pt,shorten >=1pt] (b\i.north east) -- (m\j.south west);
    \draw[-{Latex[length=2.2mm]},fmA!70,line width=1pt,dashed,shorten >=2pt] ([xshift=8pt]t\j.south) -- ([xshift=8pt]b\j.north);}
\end{tikzpicture}}

% capçalera i peu
\pagestyle{fancy}\fancyhf{}
\fancyhead[L]{\small\fmbrand}
\fancyhead[R]{\small\color{fmmuted}3r ESO · Unitat 10: Proporcionalitat. Teorema de Tales}
\fancyfoot[L]{\scriptsize\color{fmmuted}\href{https://creativecommons.org/licenses/by-nc-sa/4.0/deed.ca}{\ccbyncsa}\ \ FentMates · fentmates.org · CC BY-NC-SA 4.0}
\fancyfoot[R]{\small\color{fmmuted}\thepage}
\renewcommand{\headrulewidth}{0.4pt}\renewcommand{\headrule}{\color{fmmuted!40}\hrule height \headrulewidth}
\fancypagestyle{primera}{\fancyhf{}\renewcommand{\headrulewidth}{0pt}
  \fancyfoot[L]{\scriptsize\color{fmmuted}\href{https://creativecommons.org/licenses/by-nc-sa/4.0/deed.ca}{\ccbyncsa}\ \ FentMates · fentmates.org · CC BY-NC-SA 4.0}
  \fancyfoot[R]{\small\color{fmmuted}\thepage}}

% targeta de l'índex
\newcommand{\itile}[4]{\begin{tcolorbox}[enhanced,colback=white,colframe=fmblue!35,arc=3mm,boxrule=1pt,boxsep=0pt,left=6pt,right=6pt,top=5pt,bottom=5pt,before skip=0pt,after skip=7pt]
  \hyperlink{#1}{\begin{minipage}[c]{0.9cm}\centering{\color{fmblue}\LARGE #3}\end{minipage}\hspace{6pt}%
  \begin{minipage}[c]{0.78\linewidth}{\color{fmblue}\heavy #2}\\[-2pt]{\color{fmnavy}\semi\large #4}\end{minipage}}\end{tcolorbox}}

% taula amb vores arrodonides: el contingut es retalla amb la forma (sense @{} a les vores ni entre columnes)
\newsavebox\tabclipbox
\newcommand{\taulaclip}[3][3mm]{\sbox\tabclipbox{#3}\begin{tikzpicture}[baseline=(tcn.base)]
  \begin{scope}\clip[rounded corners=#1] (0,-\dp\tabclipbox) rectangle (\wd\tabclipbox,\ht\tabclipbox);
  \node[inner sep=0pt,outer sep=0pt,anchor=base west] (tcn) at (0,0) {\usebox\tabclipbox};\end{scope}
  \draw[#2,rounded corners=#1] (0,-\dp\tabclipbox) rectangle (\wd\tabclipbox,\ht\tabclipbox);
\end{tikzpicture}}
\newcommand{\taularod}[3][fmblue]{\taulaclip{#1,line width=1.2pt}{\renewcommand{\arraystretch}{1.5}\begin{tabular}{#2}#3\end{tabular}}}
% balança
\newcommand{\balanca}[3]{\begin{tikzpicture}[scale=#3,transform shape,every node/.style={anchor=center}]
  \fill[fmnavy] (-0.45,-1.9) rectangle (0.45,-1.75);
  \fill[fmnavy!80] (-0.12,-1.75) rectangle (0.12,0);
  \fill[fmgold] (0,0.05) circle (0.13);
  \draw[fmnavy,line width=2.4pt,line cap=round] (-2,0) -- (2,0);
  \foreach \s in {-1,1} {
    \draw[fmmuted,line width=0.8pt] (\s*2,0) -- (\s*2-0.75,-1.05) (\s*2,0) -- (\s*2+0.75,-1.05);
    \fill[fmyellow!70,draw=fmgold,line width=1pt] (\s*2-0.95,-1.05) arc[start angle=180,end angle=360,x radius=0.95,y radius=0.32] -- cycle;}
  \node[font=\Large,fill=fmA!15,draw=fmA,rounded corners=1.5mm,inner sep=3pt] at (-2,-0.68) {$#1$};
  \node[font=\Large,fill=fmB!15,draw=fmB,rounded corners=1.5mm,inner sep=3pt] at (2,-0.68) {$#2$};
\end{tikzpicture}}
% eixos amb quadrícula (dins d'un tikzpicture): xmin, xmax, ymin, ymax
\newcommand{\eixos}[4]{\draw[fmmuted!22,line width=0.4pt] (#1,#3) grid (#2,#4);
  \draw[-{Latex[length=2mm]},fmmuted,line width=0.8pt] (#1-0.2,0)--(#2+0.4,0) node[right,font=\small]{$x$};
  \draw[-{Latex[length=2mm]},fmmuted,line width=0.8pt] (0,#3-0.2)--(0,#4+0.4) node[above,font=\small]{$y$};}
% taula de traduccions (diccionari) amb les vores arrodonides
\newcommand{\dic}[2]{#1 & {\large$#2$}\\}
\newcommand{\dicratio}{0.62}
\newcommand{\dichead}{\begin{tabular}{@{}>{\raggedright}p{\dimexpr\dicratio\linewidth-12pt}>{\centering\arraybackslash}p{\dimexpr\linewidth-\dicratio\linewidth-12pt}@{}}\faComment[regular]\ Ho diem així\dots & \dots\ així\end{tabular}}
\newenvironment{dicbox}{\begin{tcolorbox}[enhanced,arc=3mm,boxrule=1.2pt,colframe=fmblue,colback=white,coltitle=white,fonttitle=\heavy,
  title={\dichead},left=0pt,right=0pt,top=0pt,bottom=1.5mm,boxsep=0pt,lefttitle=6pt,righttitle=6pt,before skip=0pt,after skip=0pt]\rowcolors{1}{white}{fmsurface}%
  \begin{tabular}{@{\hspace{6pt}}>{\raggedright}p{\dimexpr\dicratio\linewidth-12pt}>{\centering\arraybackslash}p{\dimexpr\linewidth-\dicratio\linewidth-12pt}@{\hspace{6pt}}}}
  {\end{tabular}\end{tcolorbox}}
% resum
\newcommand{\rtile}[3]{\begin{tcolorbox}[enhanced,colback=white,colframe=fmblue!45,arc=3mm,boxrule=1pt,left=3mm,right=3mm,top=2.5mm,bottom=2.5mm,fontupper=\large,
  title={\tcbox[on line,colback=fmblue,colframe=fmblue,arc=2mm,boxsep=0pt,left=4pt,right=4pt,top=1pt,bottom=1pt]{\color{white}\small\heavy #1}\ \ #2},coltitle=fmnavy,colbacktitle=fmdefbg,fonttitle=\semi,toptitle=0.5mm,bottomtitle=0.5mm]#3\end{tcolorbox}}
% portada: #1 subtítol a mà
\newcommand{\portada}[1]{\thispagestyle{primera}
\begin{tikzpicture}[remember picture,overlay]
  \fill[fmblue] (current page.north west) rectangle ([yshift=-9.4cm]current page.north east);
  \begin{scope}
    \clip (current page.north west) rectangle ([yshift=-9.4cm]current page.north east);
    \node[anchor=north west,inner sep=0pt] at ([xshift=6.6cm,yshift=-1.3cm]current page.north west) {\fmcard};
  \end{scope}
  \node[anchor=north west,inner sep=0pt] at ([xshift=1.6cm,yshift=-1.1cm]current page.north west) {\tcbox[colback=white,colframe=white,arc=2.5mm,boxsep=0pt,left=7pt,right=7pt,top=5pt,bottom=5pt]{\includegraphics[height=0.75cm]{logo-fentmates.pdf}}};
\end{tikzpicture}
\vspace*{0.55cm}
\tcbox[colback=white!15!fmblue,colframe=white!15!fmblue,arc=2.5mm,boxsep=0pt,left=6pt,right=6pt,top=3pt,bottom=3pt]{\color{white}\small\bfseries Apunts de teoria · 3r ESO · Matemàtiques}

\vspace{2pt}
{\raggedright\color{white}\fontsize{30}{36}\heavy\selectfont{\color{fmyellow}Unitat 10}\\[2pt]Proporcionalitat. Teorema de Tales\par}
\vspace{8pt}
{\color{white}\comic\Large #1\par}
\vspace*{\dimexpr 7.4cm-\pagetotal\relax}
}
\newcommand{\llegenda}{\par\vspace{16pt}
{\color{fmnavy}\heavy\Large\faPalette\ \ Com es llegeixen aquests apunts}\par\smallskip
\begin{tcbraster}[raster columns=4,raster equal height,raster column skip=8pt,raster before skip=12pt]
\tcbox[colback=fmdefbg,colframe=fmdef,arc=3mm,boxrule=1.2pt,halign=center,width=\linewidth]{{\color{fmdef}\Large\faKey}\\[2pt]{\color{fmdef}\heavy Idea clau}\\[-1pt]{\small el que cal saber}}
\tcbox[colback=fmexbg,colframe=fmex,arc=3mm,boxrule=1.2pt,halign=center,width=\linewidth]{{\color{fmex}\Large\faPencil*}\\[2pt]{\color{fmex}\heavy Exemple}\\[-1pt]{\small pas a pas}}
\tcbox[colback=fmsf!6,colframe=fmsf,arc=3mm,boxrule=1.2pt,halign=center,width=\linewidth]{{\color{fmsf}\Large\faExclamationTriangle}\\[2pt]{\color{fmsf}\heavy Compte!}\\[-1pt]{\small errors típics}}
\tcbox[colback=fmyellow!20,colframe=fmgold,arc=3mm,boxrule=1.2pt,halign=center,width=\linewidth]{{\color{fmgold!80!black}\Large\faLightbulb[regular]}\\[2pt]{\color{fmnavy}\heavy Truc}\\[-1pt]{\small per a recordar-ho}}
\end{tcbraster}}
\newcommand{\resumtitol}{\clearpage
\hypertarget{resum}{}
{\color{fmnavy}\fontsize{23}{27}\heavy\selectfont\faClipboardList\ \ En resum}\par\vspace{5pt}{\color{fmblue}\titlerule[2pt]}
\frase{Tota la unitat en una pàgina.}}
% els 5 passos per a resoldre problemes: #2 = text del pas 2
\newcommand{\passosproblema}[2]{\begin{center}
\begin{tikzpicture}[s/.style={draw=fmblue,fill=fmdefbg,rounded corners=3mm,line width=1.2pt,minimum width=2.75cm,minimum height=2.45cm,align=center,font=\small},
  ic/.style={font=\LARGE,text=fmblue}]
  \foreach \i/\ic/\t in {1/\faBookOpen/{\textbf{Llegim}\\i entenem},2/\faQuestionCircle/{#1},3/\faEquals/{#2},4/\faCogs/{\textbf{Resolem}},5/\faCheckDouble/{\textbf{Comprovem}\\i responem}} {
    \node[s] (s\i) at (3.3*\i,0) {};
    \node[ic] at ([yshift=-0.5cm]s\i.north) {\ic};
    \node[align=center,font=\small,anchor=north] at ([yshift=-0.95cm]s\i.north) {\t};
    \node[circle,fill=fmblue,text=white,font=\small\heavy,inner sep=1.5pt,minimum size=1.3em] at (s\i.north west) {\i};}
  \foreach \i [evaluate=\i as \j using int(\i+1)] in {1,...,4} \draw[-{Latex[length=3mm]},fmmuted,line width=1.5pt] (s\i) -- (s\j);
\end{tikzpicture}
\end{center}}

\newcommand{\lo}{\mathopen{]}}\newcommand{\ro}{\mathclose{[}}
% punts ple i buit, nombres dels eixos i taula de valors
\newcommand{\pple}[2][fmblue]{\fill[#1] #2 circle (3.6pt);}
\newcommand{\nums}[4]{\foreach \k in {#1,...,#2} {\ifnum\k=0\else\node[below,font=\scriptsize,text=fmmuted] at (\k,0) {$\k$};\fi}
  \foreach \k in {#3,...,#4} {\ifnum\k=0\else\node[left,font=\scriptsize,text=fmmuted] at (0,\k) {$\k$};\fi}}
\newcommand{\tv}[3][fmblue]{\taulaclip{#1,line width=1.2pt}{\renewcommand{\arraystretch}{1.45}%
  \begin{tabular}{>{\columncolor{#1}\color{white}\heavy}c*{12}{w{c}{0.8cm}}}#2\\\rowcolor{#1!8}\cellcolor{#1}#3\end{tabular}}}
% etiqueta d'una recta o corba
\newcommand{\etq}[3]{\node[fill=white,inner sep=1.5pt,text=#1,font=\small] at #2 {#3};}

\begin{document}
\portada{Segments proporcionals, figures semblants i com mesurar el que no podem tocar.}

{\color{fmnavy}\heavy\Large\faMapSigns\ \ Què aprendràs?}\par\smallskip
\begin{minipage}[t]{0.485\linewidth}
\itile{sec1}{10.1}{\faRulerHorizontal}{Raons i proporcions}
\itile{sec2}{10.2}{\faGripLines}{El teorema de Tales}
\itile{sec3}{10.3}{\faCut}{Dividir un segment}
\itile{sec4}{10.4}{\faIcon{play}}{Triangles en posició de Tales}
\end{minipage}\hfill
\begin{minipage}[t]{0.485\linewidth}
\itile{sec5}{10.5}{\faClone}{Figures semblants}
\itile{sec6}{10.6}{\faBalanceScale}{Criteris de semblança}
\itile{sec7}{10.7}{\faSun}{Altures i distàncies}
\itile{sec8}{10.8}{\faMap}{Escales, àrees i volums}
\itile{resum}{\phantom{10.8}}{\faClipboardList}{En resum}
\end{minipage}
\llegenda

% ============================================================ 10.1
\clearpage
\seccio{sec1}{Raons i proporcions entre segments}{Comparem longituds dividint-les.}

\vfill
\begin{minipage}[c]{0.45\linewidth}\centering
\begin{tikzpicture}[x=0.75cm]
  \draw[fmA,line width=3pt] (0,1) -- (6,1); \foreach \k in {0,...,6} \draw[fmA,line width=0.8pt] (\k,0.88) -- (\k,1.12);
  \node[left,font=\small\heavy,text=fmA] at (0,1) {$AB$}; \node[right,font=\small,text=fmA] at (6,1) {6 cm};
  \draw[fmB,line width=3pt] (0,0) -- (4,0); \foreach \k in {0,...,4} \draw[fmB,line width=0.8pt] (\k,-0.12) -- (\k,0.12);
  \node[left,font=\small\heavy,text=fmB] at (0,0) {$CD$}; \node[right,font=\small,text=fmB] at (4,0) {4 cm};
\end{tikzpicture}\\[6pt]
{\Large$\dfrac{\overline{AB}}{\overline{CD}}=\dfrac64=1{,}5$}\\[2pt]
{\comic\color{fmmuted}$AB$ és una vegada i mitja $CD$}
\end{minipage}\hfill
\begin{minipage}[c]{0.52\linewidth}
\begin{idea}[Raó i proporció]
La \textbf{raó} de dos segments és el quocient de les seues longituds, \textbf{en la mateixa unitat}. No té unitats.\\[4pt]
Quatre segments són \textbf{proporcionals} si $\dfrac ab=\dfrac cd$: és una \textbf{proporció}.\\[4pt]
\textbf{Propietat fonamental}: \ $\dfrac ab=\dfrac cd\ \Longleftrightarrow\ a\cdot d=b\cdot c$
\end{idea}
\end{minipage}

\vfill
\begin{tcbraster}[raster columns=2,raster equal height,raster column skip=10pt]
\begin{tcolorbox}[fmbase,colback=fmK!8,colframe=fmK,coltitle=white,title={Són proporcionals?}]
$3,\,5,\,6,\,10$: \ $3\cdot10=5\cdot6=30$ \ \be{} \ sí\\[4pt]
$2,\,3,\,4,\,5$: \ $2\cdot5=10\neq3\cdot4=12$ \ \mal{\text{no}}
\end{tcolorbox}
\begin{tcolorbox}[fmbase,colback=fmA!8,colframe=fmA,coltitle=white,title={El quart proporcional}]
$x$ amb $\dfrac ab=\dfrac cx$: \ multipliquem \textbf{en creu}:\par\smallskip
{\centering$\dfrac46=\dfrac{10}{x}\ \to\ 4x=60\ \to\ \boxed{x=15}$\par}
\end{tcolorbox}
\end{tcbraster}

\vfill
\begin{exemple}[Exemple: \ $\overline{AB}=0{,}6$ m, $\overline{CD}=40$ cm. Quant fa $\overline{GH}$ si $\frac{\overline{AB}}{\overline{CD}}=\frac{\overline{EF}}{\overline{GH}}$ i $\overline{EF}=9$ cm?]
\pas{1} La mateixa unitat: \ $0{,}6$ m $=60$ cm \ $\to$ \ $\dfrac{60}{40}=\dfrac32$\\[4pt]
\pas{2} $\dfrac32=\dfrac{9}{\overline{GH}}\ \to\ \overline{GH}=\dfrac{2\cdot9}{3}=\boxed{6\text{ cm}}$
\end{exemple}

\vfill
\begin{minipage}[t]{0.5\linewidth}
\begin{exemple}[Exemple: \ mitjà proporcional de 4 i 9]
$\dfrac4x=\dfrac x9\ \to\ x^2=36\ \to\ \boxed{x=6\text{ cm}}$
\end{exemple}
\end{minipage}\hfill
\begin{minipage}[t]{0.47\linewidth}
\begin{truc}[On siga la $x$]
La incògnita pot estar en qualsevol lloc. Fes sempre el mateix: \textbf{multiplica en creu} i aïlla la $x$.
\end{truc}
\end{minipage}

\vfill
% ============================================================ 10.2
\clearpage
\seccio{sec2}{El teorema de Tales}{Les paral·leles tallen les secants en trossos proporcionals.}

\vfill
\begin{minipage}[c]{0.45\linewidth}\centering
\begin{tikzpicture}[scale=0.62]
  \draw[fmnavy,line width=1.3pt] (0,0) -- (9,1.35) node[right,font=\small] {$a$};
  \draw[fmnavy,line width=1.3pt] (0,0) -- (7.1,6.39) node[above right,font=\small] {$b$};
  \foreach \p/\q in {{(2,0.3)}/{(1.601,1.441)},{(4.5,0.675)}/{(3.602,3.242)},{(8,1.2)}/{(6.403,5.764)}} {
    \draw[fmV,line width=1.3pt,shorten <=-0.5cm,shorten >=-0.5cm] \p -- \q;}
  \draw[fmA,line width=3.2pt] (2,0.3) -- (4.5,0.675); \draw[fmB,line width=3.2pt] (4.5,0.675) -- (8,1.2);
  \draw[fmA,line width=3.2pt] (1.601,1.441) -- (3.602,3.242); \draw[fmB,line width=3.2pt] (3.602,3.242) -- (6.403,5.764);
  \foreach \p/\n in {{(2,0.3)}/A,{(4.5,0.675)}/B,{(8,1.2)}/C} \fill[fmnavy] \p circle (2.6pt) node[below,font=\small] {$\n$};
  \foreach \p/\n in {{(1.601,1.441)}/{A'},{(3.602,3.242)}/{B'},{(6.403,5.764)}/{C'}} \fill[fmnavy] \p circle (2.6pt) node[left,font=\small] {$\n$};
\end{tikzpicture}\\
{\comic\small\color{fmV}les rectes morades són paral·leles}
\end{minipage}\hfill
\begin{minipage}[c]{0.52\linewidth}
\begin{idea}[Teorema de Tales]
Si unes rectes \textbf{paral·leles} tallen dues rectes secants, els segments d'una són \textbf{proporcionals} als segments corresponents de l'altra:\\[5pt]
{\centering\Large$\dfrac{\cA{\overline{AB}}}{\cB{\overline{BC}}}=\dfrac{\cA{\overline{A'B'}}}{\cB{\overline{B'C'}}}$\par}\medskip
o bé \ $\dfrac{\overline{AB}}{\overline{A'B'}}=\dfrac{\overline{BC}}{\overline{B'C'}}=\dfrac{\overline{AC}}{\overline{A'C'}}$
\end{idea}
\end{minipage}

\vfill
\begin{compte}[Els segments de les secants, no els de les paral·leles]
El teorema parla dels trossos \textbf{sobre les secants} $a$ i $b$. Els segments de les paral·leles (entre $a$ i $b$) no compleixen, en general, aquesta proporció.
\end{compte}

\vfill
\begin{tcbraster}[raster columns=2,raster equal height,raster column skip=10pt]
\begin{tcolorbox}[fmbase,colback=fmK!8,colframe=fmK,coltitle=white,title={Calcular un segment}]
$\overline{AB}=4$, $\overline{BC}=6$, $\overline{A'B'}=5$:\par\smallskip
{\centering$\dfrac46=\dfrac5x\ \to\ 4x=30\ \to\ \boxed{x=7{,}5}$\par}
\end{tcolorbox}
\begin{tcolorbox}[fmbase,colback=fmdefbg,colframe=fmdef,coltitle=white,title={El recíproc: són paral·leles?}]
Si els segments són proporcionals, les rectes \textbf{són paral·leles}:\par\smallskip
{\centering$3,\,4$ \ i \ $4{,}5,\,6$: \ $\dfrac34=0{,}75=\dfrac{4{,}5}{6}$ \ \be{}\par}
\end{tcolorbox}
\end{tcbraster}

\vfill
\begin{exemple}[Exemple: \ tres parcel·les amb partions paral·leles: 20, 30 i 40 m al carrer Major; 108 m en total al carrer Nou]
\begin{minipage}[c]{0.6\linewidth}
La raó entre els dos carrers és sempre la mateixa:\\[3pt]
$k=\dfrac{108}{20+30+40}=\dfrac{108}{90}=1{,}2$\\[4pt]
$20\cdot1{,}2=\boxed{24\text{ m}}$, \ $30\cdot1{,}2=\boxed{36\text{ m}}$, \ $40\cdot1{,}2=\boxed{48\text{ m}}$
\end{minipage}\hfill
\begin{minipage}[c]{0.36\linewidth}\centering
\begin{tikzpicture}[scale=0.4]
  \draw[fmmuted,line width=1.2pt] (0,0) -- (9,0) node[right,font=\scriptsize] {Major};
  \draw[fmmuted,line width=1.2pt] (0,3) -- (10.8,3) node[right,font=\scriptsize] {Nou};
  \foreach \x/\y in {0/0,2/2.4,5/6,9/10.8} \draw[fmex,line width=1.2pt] (\x,0) -- (\y,3);
  \node[below,font=\scriptsize] at (1,0) {20}; \node[below,font=\scriptsize] at (3.5,0) {30}; \node[below,font=\scriptsize] at (7,0) {40};
\end{tikzpicture}
\end{minipage}
\end{exemple}

\vfill
% ============================================================ 10.3
\clearpage
\seccio{sec3}{Dividir un segment}{Amb regle i compàs, i sense mesurar.}

\vfill
\begin{minipage}[c]{0.5\linewidth}\centering
\begin{tikzpicture}[scale=0.85]
  \draw[fmnavy,line width=2pt] (0,0) -- (6,0);
  \fill[fmnavy] (0,0) circle (2pt) node[below,font=\small] {$A$}; \fill[fmnavy] (6,0) circle (2pt) node[below,font=\small] {$B$};
  \draw[fmmuted,line width=1pt,-{Latex[length=2mm]}] (0,0) -- (35:6.3);
  \foreach \k in {1,...,5} \fill[fmB] (35:\k) circle (2.2pt);
  \foreach \k in {1,...,5} \node[above left,font=\scriptsize,text=fmB] at (35:\k) {\k};
  \draw[fmR,line width=1.3pt] (35:5) -- (6,0);
  \foreach \k in {1,...,4} \draw[fmR,line width=1pt,dashed] (35:\k) -- ({6*\k/5},0);
  \foreach \k in {1,...,4} \fill[fmK] ({6*\k/5},0) circle (2.2pt);
\end{tikzpicture}\\
{\comic\small\color{fmmuted}$AB$ dividit en 5 parts iguals}
\end{minipage}\hfill
\begin{minipage}[c]{0.47\linewidth}\raggedright
\pas{1} Des de $A$ traça una \textbf{semirecta} auxiliar.\\[5pt]
\pas{2} Amb el compàs, marca-hi \textbf{$n$ segments iguals}.\\[5pt]
\pas{3} Uneix l'última marca amb $B$.\\[5pt]
\pas{4} Per les altres marques, traça \textbf{paral·leles}: tallen $AB$ en $n$ parts iguals.\\[4pt]
\notab{Per Tales: segments iguals en la semirecta, segments iguals en $AB$.}
\end{minipage}

\vfill
\begin{tcolorbox}[fmbase,colback=fmdefbg,colframe=fmdef,coltitle=white,title={En parts proporcionals a 2, 3 i 4}]
\begin{minipage}[c]{0.55\linewidth}
\textbf{Construcció}: marquem $2+3+4=9$ segments iguals i tracem paral·leles només per les marques 2 i $2+3=5$.\\[5pt]
\textbf{Càlcul}: és un \textbf{repartiment proporcional}:\\[3pt]
$\dfrac{18}{2+3+4}=2$ \ $\to$ \ $4$ cm, \ $6$ cm, \ $8$ cm
\end{minipage}\hfill
\begin{minipage}[c]{0.42\linewidth}\centering
\begin{tikzpicture}[x=0.35cm,y=0.35cm]
  \draw[fmA,line width=4pt] (0,0) -- (4,0); \draw[fmB,line width=4pt] (4,0) -- (10,0); \draw[fmK,line width=4pt] (10,0) -- (18,0);
  \node[below=3pt,font=\small] at (2,0) {4}; \node[below=3pt,font=\small] at (7,0) {6}; \node[below=3pt,font=\small] at (14,0) {8};
  \node[above=3pt,font=\small] at (9,0) {18 cm};
\end{tikzpicture}
\end{minipage}
\end{tcolorbox}

\vfill
\begin{exemple}[Exemples]
\begin{minipage}[t]{0.48\linewidth}
15 cm en 6 parts iguals: \ $15:6=\boxed{2{,}5\text{ cm}}$\\[4pt]
20 cm proporcional a 1, 3 i 4:\\ $20:8=2{,}5$ \ $\to$ \ \boxed{2{,}5}, \boxed{7{,}5}, \boxed{10} cm
\end{minipage}\hfill
\begin{minipage}[t]{0.49\linewidth}
$P$ divideix 28 cm amb $\frac{\overline{AP}}{\overline{PB}}=\frac34$:\\[2pt]
parts proporcionals a 3 i 4: \ $28:7=4$\\[2pt]
$\overline{AP}=\boxed{12\text{ cm}}$, \ $\overline{PB}=\boxed{16\text{ cm}}$
\end{minipage}
\end{exemple}

\vfill
% ============================================================ 10.4
\clearpage
\seccio{sec4}{Triangles en posició de Tales}{Un angle comú i els costats de davant paral·lels.}

\vfill
\begin{minipage}[c]{0.45\linewidth}\centering
\begin{tikzpicture}[scale=0.75]
  \fill[fmB!15] (0,0) -- (7,0) -- (5,4.2) -- cycle;
  \fill[fmA!25] (0,0) -- (2.8,0) -- (2,1.68) -- cycle;
  \draw[fmnavy,line width=1.4pt] (0,0) -- (7,0) -- (5,4.2) -- cycle;
  \draw[fmA,line width=2pt] (2.8,0) -- (2,1.68);
  \draw[fmB,line width=2pt] (7,0) -- (5,4.2);
  \node[below left,font=\small] at (0,0) {$A$};
  \node[below,font=\small] at (2.8,0) {$B'$}; \node[below,font=\small] at (7,0) {$B$};
  \node[above left,font=\small] at (2,1.68) {$C'$}; \node[above,font=\small] at (5,4.2) {$C$};
\end{tikzpicture}\\
{\comic\small\color{fmmuted}$B'C'\parallel BC$}
\end{minipage}\hfill
\begin{minipage}[c]{0.52\linewidth}
\begin{idea}[Posició de Tales]
Els triangles $AB'C'$ i $ABC$ tenen l'angle $\hat A$ \textbf{comú} i els costats de davant \textbf{paral·lels}. Aleshores tots els costats són proporcionals:\\[5pt]
{\centering\Large$\dfrac{\overline{AB'}}{\overline{AB}}=\dfrac{\overline{AC'}}{\overline{AC}}=\dfrac{\cA{\overline{B'C'}}}{\cB{\overline{BC}}}$\par}
\end{idea}
\end{minipage}

\vfill
\begin{compte}[Costats sencers, des de $A$]
Per al costat paral·lel, compara costats \textbf{sencers}: \quad \be{\dfrac{\overline{AB'}}{\overline{AB}}=\dfrac{\overline{B'C'}}{\overline{BC}}}\qquad \mal{\dfrac{\overline{AB'}}{\overline{B'B}}=\dfrac{\overline{B'C'}}{\overline{BC}}}
\end{compte}

\vfill
\begin{exemple}[Exemple: \ $\overline{AB'}=4$, $\overline{B'B}=6$, $\overline{BC}=12{,}5$. Quant fa $\overline{B'C'}$?]
$\overline{AB}=4+6=10$ \qquad $\dfrac{4}{10}=\dfrac{x}{12{,}5}\ \to\ x=\dfrac{4\cdot12{,}5}{10}=\boxed{5}$
\end{exemple}

\vfill
\begin{exemple}[Exemple: \ l'amplària d'un riu]
\begin{minipage}[c]{0.4\linewidth}\centering
\begin{tikzpicture}[scale=0.15]
  \fill[fmA!18] (-2,0.6) rectangle (22,15); \node[font=\scriptsize,text=fmA] at (16,11) {riu};
  \draw[fmmuted,line width=1pt] (-2,0) -- (22,0);
  \draw[fmnavy,line width=1.2pt] (0,15) -- (20,-10); \draw[fmnavy,line width=1.2pt] (0,0) -- (20,0);
  \draw[fmex,line width=2pt] (0,0) -- (0,15); \draw[fmex,line width=2pt] (20,0) -- (20,-10);
  \foreach \p/\n/\o in {{(0,15)}/A/left,{(0,0)}/B/below left,{(12,0)}/C/above,{(20,0)}/D/above right,{(20,-10)}/E/right} \fill[fmnavy] \p circle (8pt) node[\o,font=\scriptsize] {$\n$};
  \node[left,font=\scriptsize\heavy] at (0,7.5) {$x$}; \node[below,font=\scriptsize] at (6,0) {12}; \node[below,font=\scriptsize] at (16,0) {8}; \node[right,font=\scriptsize] at (20,-5) {10};
\end{tikzpicture}
\end{minipage}\hfill
\begin{minipage}[c]{0.57\linewidth}
Clavem estaques: els triangles tenen l'angle en $C$ igual (oposats pel vèrtex) i els costats $AB$ i $DE$ paral·lels: posició de Tales \textbf{girada}.\\[4pt]
$\dfrac{x}{10}=\dfrac{12}{8}\ \to\ x=\dfrac{10\cdot12}{8}=\boxed{15\text{ m}}$
\end{minipage}
\end{exemple}

\vfill
\begin{truc}[Tot és proporcional]
Les altures, els perímetres\dots\ dels dos triangles també estan en la mateixa raó.
\end{truc}

\vfill
% ============================================================ 10.5
\clearpage
\seccio{sec5}{Figures semblants. Raó de semblança}{Mateixa forma, diferent grandària.}

\vfill
\begin{minipage}[c]{0.45\linewidth}\centering
\begin{tikzpicture}[scale=0.5]
  \filldraw[fill=fmA!20,draw=fmA,line width=1.4pt] (0,0) -- (4,0) -- (4,1.5) -- (1.5,1.5) -- (1.5,3) -- (0,3) -- cycle;
  \node[below,font=\small,text=fmA] at (2,0) {4};
  \filldraw[fill=fmB!20,draw=fmB,line width=1.4pt] (5,0) -- (11,0) -- (11,2.25) -- (7.25,2.25) -- (7.25,4.5) -- (5,4.5) -- cycle;
  \node[below,font=\small,text=fmB] at (8,0) {6};
\end{tikzpicture}\\[3pt]
{\Large$k=\dfrac64=1{,}5$}
\end{minipage}\hfill
\begin{minipage}[c]{0.52\linewidth}
\begin{idea}[Figures semblants]
Una és una \textbf{ampliació} o una \textbf{reducció} de l'altra:\\[3pt]
\faCheck\ angles corresponents \textbf{iguals};\\[2pt]
\faCheck\ costats homòlegs \textbf{proporcionals}.\\[4pt]
\textbf{Raó de semblança}: \ $k=\dfrac{a'}{a}=\dfrac{b'}{b}=\dots$\\[3pt]
\chip{fmK}{$k>1$} ampliació \ \chip{fmblue}{$k=1$} igual \ \chip{fmR}{$k<1$} reducció
\end{idea}
\end{minipage}

\vfill
\begin{compte}[Calen les dues condicions]
Un quadrat i un rectangle tenen els angles iguals però \textbf{no} són semblants; un quadrat i un rombe tenen els costats proporcionals però \textbf{tampoc}.
\end{compte}

\vfill
\begin{tcolorbox}[enhanced,colback=fmyellow!18,colframe=fmgold,arc=4mm,boxrule=1.6pt,before upper=\centering,top=2mm,bottom=2mm]
{\Large Longituds: $\times\,k$\qquad\qquad Àrees: $\times\,k^2$}\par\smallskip
{\small si dupliquem els costats ($k=2$), l'àrea es multiplica per 4: \ $kb\cdot kh=k^2\cdot bh$}
\end{tcolorbox}

\vfill
\begin{exemple}[Exemples]
\begin{minipage}[t]{0.48\linewidth}
Rectangle $4\times6$ i un semblant amb el costat menut de 10:\\
$k=\frac{10}{4}=2{,}5$, \ el gran: $6\cdot2{,}5=\boxed{15\text{ cm}}$\\[5pt]
Són semblants $3\times5$ i $6\times9$?\\
$\frac63=2\neq\frac95=1{,}8$ \ \mal{\text{no}}
\end{minipage}\hfill
\begin{minipage}[t]{0.49\linewidth}
Triangle de 15 cm de perímetre i 12 cm$^2$, $k=3$:\\
$P'=15\cdot3=\boxed{45\text{ cm}}$, \ $A'=12\cdot9=\boxed{108\text{ cm}^2}$\\[5pt]
Foto de $10\times15$ amb l'àrea \textbf{doble}:\\
$k^2=2\to k=\sqrt2$: \ \boxed{14{,}14\times21{,}21\text{ cm}}
\end{minipage}
\end{exemple}

\vfill
\begin{truc}[Doblar l'àrea no és doblar els costats]
Si dobles els costats, l'àrea es fa \textbf{4 vegades} més gran. Per a doblar l'àrea, multiplica els costats per $\sqrt2\approx1{,}41$.
\end{truc}

\vfill
% ============================================================ 10.6
\clearpage
\seccio{sec6}{Criteris de semblança de triangles}{No cal comprovar-ho tot: amb un criteri n'hi ha prou.}

\vfill
\begin{tcbraster}[raster columns=3,raster equal height,raster column skip=8pt]
\begin{tcolorbox}[fmbase,colback=fmA!8,colframe=fmA,coltitle=white,title={1r · Angles},before upper=\centering]
\begin{tikzpicture}[scale=0.42]
  \draw[fmA,line width=1.2pt] (0,0) -- (3,0) -- (1,2) -- cycle; \draw[fmA,line width=1.2pt] (3.6,0) -- (8.1,0) -- (5.1,3) -- cycle;
  \draw[fmR] (0.5,0) arc (0:45:0.5); \draw[fmR] (4.1,0) arc (0:45:0.5);
  \draw[fmK] (2.55,0) arc (180:135:0.45); \draw[fmK] (7.65,0) arc (180:135:0.45);
\end{tikzpicture}\par\smallskip
\textbf{Dos angles} iguals\par{\small (el tercer també)}
\end{tcolorbox}
\begin{tcolorbox}[fmbase,colback=fmB!8,colframe=fmB,coltitle=white,title={2n · Costats},before upper=\centering]
\begin{tikzpicture}[scale=0.42]
  \draw[fmB,line width=1.2pt] (0,0) -- (3,0) -- (1,2) -- cycle; \draw[fmB,line width=1.2pt] (3.6,0) -- (8.1,0) -- (5.1,3) -- cycle;
\end{tikzpicture}\par\smallskip
Els \textbf{tres costats} proporcionals\par{\small$\frac{a'}{a}=\frac{b'}{b}=\frac{c'}{c}$}
\end{tcolorbox}
\begin{tcolorbox}[fmbase,colback=fmK!8,colframe=fmK,coltitle=white,title={3r · Costats i angle},before upper=\centering]
\begin{tikzpicture}[scale=0.42]
  \draw[fmK,line width=1.2pt] (0,0) -- (3,0) -- (1,2) -- cycle; \draw[fmK,line width=1.2pt] (3.6,0) -- (8.1,0) -- (5.1,3) -- cycle;
  \draw[fmR,line width=1pt] (0.5,0) arc (0:63:0.5); \draw[fmR,line width=1pt] (4.1,0) arc (0:63:0.5);
\end{tikzpicture}\par\smallskip
\textbf{Dos costats} proporcionals i l'\textbf{angle} entre ells igual
\end{tcolorbox}
\end{tcbraster}

\vfill
\begin{truc}[Comprovar-ho amb els costats]
Ordena els costats de menor a major i divideix-los en aquest ordre:\\[3pt]
$(4,\,6,\,7)$ i $(6,\,9,\,10{,}5)$: \ $\dfrac64=\dfrac96=\dfrac{10{,}5}{7}=1{,}5$ \ $\to$ \ semblants, $k=1{,}5$
\end{truc}

\vfill
{\large Dos triangles en \textbf{posició de Tales} sempre són semblants. Dos triangles \textbf{rectangles} ho són si tenen un angle agut igual.}

\vfill
\begin{exemple}[Exemples]
\begin{minipage}[t]{0.48\linewidth}
$5,\,7,\,8$ i $15,\,21,\,24$: \ $\frac{15}{5}=\frac{21}{7}=\frac{24}{8}=3$\\ \be{} semblants, $\boxed{k=3}$\\[5pt]
Angles $50^\circ,\,65^\circ$ i $65^\circ,\,65^\circ$: el tercer és $65^\circ$ i $50^\circ$: \be{} \textbf{mateixos angles}
\end{minipage}\hfill
\begin{minipage}[t]{0.49\linewidth}
$\hat A=\hat A'=40^\circ$, \ costats $6,\,9$ i $8,\,12$:\\
$\frac86=\frac{12}{9}=\frac43$ \ \be{} criteri 3r\\[3pt]
Si $\overline{BC}=6{,}1$: \ $\overline{B'C'}=6{,}1\cdot\frac43\approx\boxed{8{,}13}$
\end{minipage}
\end{exemple}

\vfill
\begin{exemple}[Exemple: \ altura sobre la hipotenusa d'un triangle rectangle de catets 15 i 20]
\begin{minipage}[c]{0.6\linewidth}
Hipotenusa: $\sqrt{15^2+20^2}=25$.\\[3pt]
$ABC$ i $HBA$ són rectangles i comparteixen $\hat B$: \textbf{semblants}.\\[3pt]
$\dfrac{h}{20}=\dfrac{15}{25}\ \to\ h=\dfrac{15\cdot20}{25}=\boxed{12\text{ cm}}$
\end{minipage}\hfill
\begin{minipage}[c]{0.36\linewidth}\centering
\begin{tikzpicture}[scale=0.17]
  \draw[fmex,line width=1.3pt] (0,0) -- (25,0) -- (9,12) -- cycle;
  \draw[fmK,line width=1.3pt,dashed] (9,12) -- (9,0);
  \node[below,font=\scriptsize] at (9,0) {$H$}; \node[above,font=\scriptsize] at (9,12) {$A$}; \node[below left,font=\scriptsize] at (0,0) {$B$}; \node[below right,font=\scriptsize] at (25,0) {$C$};
  \node[left,font=\scriptsize] at (4,6.5) {15}; \node[right,font=\scriptsize] at (17,6.5) {20}; \node[right,font=\scriptsize\heavy,text=fmK!80!black] at (9,5) {$h$};
\end{tikzpicture}
\end{minipage}
\end{exemple}

\vfill
% ============================================================ 10.7
\clearpage
\seccio{sec7}{Aplicacions: altures i distàncies}{Busca dos triangles semblants i escriu la proporció.}

\vfill
\begin{minipage}[c]{0.48\linewidth}\centering
\begin{tikzpicture}[scale=0.42]
  \fill[fmyellow!40] (-0.5,-0.4) rectangle (16,0);
  \draw[fmmuted,line width=1pt] (-0.5,0) -- (16,0);
  \draw[fmK,line width=2.5pt] (0,0) -- (0,1.5); \draw[fmmuted,line width=2.5pt] (0,0) -- (2,0);
  \draw[fmgold,dashed,line width=1pt] (0,1.5) -- (2,0);
  \fill[fmB!40] (4,0) rectangle (5.5,9); \draw[fmB,line width=1.2pt] (4,0) rectangle (5.5,9);
  \draw[fmmuted,line width=2.5pt] (5.5,0) -- (15.5,0);
  \draw[fmgold,dashed,line width=1pt] (5.5,9) -- (15.5,0);
  \node[left,font=\scriptsize] at (0,0.75) {1,5}; \node[below,font=\scriptsize] at (1,-0.4) {2};
  \node[left,font=\scriptsize\heavy] at (4,4.5) {$H$}; \node[below,font=\scriptsize] at (10.5,-0.4) {12};
  \node[text=fmgold,font=\Large] at (14,7) {\faSun};
\end{tikzpicture}\\
{\comic\small\color{fmmuted}rajos de Sol paral·lels: triangles semblants}
\end{minipage}\hfill
\begin{minipage}[c]{0.49\linewidth}
\begin{idea}[Altures amb ombres]
Al mateix moment, l'altura i l'ombra són \textbf{proporcionals}:\\[4pt]
{\centering\Large$\dfrac{h_{\text{pal}}}{s_{\text{pal}}}=\dfrac{H}{S}$\par}\smallskip
$\dfrac{1{,}5}{2}=\dfrac{H}{12}\ \to\ H=\dfrac{1{,}5\cdot12}{2}=\boxed{9\text{ m}}$
\end{idea}
\end{minipage}

\vfill
\begin{minipage}[c]{0.48\linewidth}\centering
\begin{tikzpicture}[scale=0.42]
  \draw[fmmuted,line width=1pt] (-0.5,0) -- (12,0);
  \draw[fmA,line width=2.5pt] (0,0) -- (0,1.6); \fill[fmA] (0,1.6) circle (4pt);
  \draw[fmK!80!black,line width=3pt] (11,0) -- (11,7.2); \fill[fmK] (11,7.2) circle (9pt);
  \fill[fmblue] (1.8,-0.12) rectangle (2.2,0.05);
  \draw[fmgold,line width=1pt] (0,1.6) -- (2,0) -- (11,7.2);
  \node[left,font=\scriptsize] at (0,0.8) {1,6}; \node[below,font=\scriptsize] at (1,0) {2}; \node[below,font=\scriptsize] at (6.5,0) {9};
  \node[right,font=\scriptsize\heavy] at (11,3.6) {$H$}; \node[above,font=\scriptsize,text=fmblue] at (2,0.1) {mirall};
\end{tikzpicture}
\end{minipage}\hfill
\begin{minipage}[c]{0.49\linewidth}
\begin{tcolorbox}[fmbase,colback=fmdefbg,colframe=fmdef,coltitle=white,title={\faIcon{eye}\ \ Amb un mirall}]
El raig es reflecteix amb el mateix angle: els dos triangles rectangles són \textbf{semblants}.\par\smallskip
{\centering$\dfrac{H}{9}=\dfrac{1{,}6}{2}\ \to\ H=\boxed{7{,}2\text{ m}}$\par}
\end{tcolorbox}
\end{minipage}

\vfill
\begin{exemple}[Exemple: \ un arbre fa 6 m d'ombra i una persona d'1,7 m, 1,2 m]
$\dfrac{H}{6}=\dfrac{1{,}7}{1{,}2}\ \to\ H=\dfrac{6\cdot1{,}7}{1{,}2}=\boxed{8{,}5\text{ m}}$
\end{exemple}

\vfill
\begin{exemple}[Exemple: \ una persona d'1,8 m està a 4 m d'un fanal de 6 m. Quant fa la seua ombra?]
Ara la llum ix d'\textbf{un punt}: triangles en posició de Tales. Si l'ombra fa $s$:\\[3pt]
$\dfrac{1{,}8}{s}=\dfrac{6}{s+4}\ \to\ 1{,}8\,(s+4)=6s\ \to\ 7{,}2=4{,}2s\ \to\ s\approx\boxed{1{,}71\text{ m}}$
\end{exemple}

\vfill
\begin{truc}[Tales i la piràmide]
Diuen que Tales de Milet va calcular l'altura de la gran piràmide esperant el moment en què la seua ombra feia el \textbf{mateix} que la seua altura.
\end{truc}

\vfill
% ============================================================ 10.8
\clearpage
\seccio{sec8}{Escales. Àrees i volums}{Plànols, mapes i maquetes són figures semblants a la realitat.}

\vfill
\begin{minipage}[c]{0.55\linewidth}
\begin{idea}[Escala]
{\centering\large$\text{escala}=\dfrac{\text{mesura en el plànol}}{\text{mesura real}}$\par}\smallskip
\textbf{Escala $1:n$}: 1 unitat del plànol són $n$ unitats reals.\\[3pt]
\chip{fmA}{real $=$ plànol $\times\,n$} \ \chip{fmB}{plànol $=$ real $:\,n$}
\end{idea}
\end{minipage}\hfill
\begin{minipage}[c]{0.42\linewidth}\centering
\begin{tikzpicture}[x=0.9cm]
  \foreach \k in {0,...,4} {\ifodd\k \fill[white,draw=fmnavy] (\k,0) rectangle (\k+1,0.3); \else \fill[fmnavy] (\k,0) rectangle (\k+1,0.3); \fi}
  \foreach \k/\l in {0/0,1/0,5/2,5/2} {}
  \node[below,font=\scriptsize] at (0,0) {0}; \node[below,font=\scriptsize] at (5,0) {2 km};
\end{tikzpicture}\\[2pt]
{\comic\small\color{fmmuted}escala gràfica: el segment sencer són 2 km}\\[6pt]
{\large$1:50\,000$}\\
{\small 1 cm $\to$ 50\,000 cm $=$ 500 m}
\end{minipage}

\vfill
\begin{exemple}[Exemples]
\begin{minipage}[t]{0.48\linewidth}
Mapa 1:200\,000, 7,5 cm: \ $7{,}5\cdot200\,000=1\,500\,000$ cm $=\boxed{15\text{ km}}$\\[5pt]
Habitació de 4,2 m $\times$ 3,5 m a 1:50:\\ $420:50=\boxed{8{,}4\text{ cm}}$, \ $350:50=\boxed{7\text{ cm}}$
\end{minipage}\hfill
\begin{minipage}[t]{0.49\linewidth}
12 km són 4 cm en el mapa:\\ $\dfrac{4}{1\,200\,000}=\dfrac{1}{300\,000}$ \ $\to$ \ \boxed{1:300\,000}\\[4pt]
Un llac de 3 cm$^2$: 1 cm són 3 km, 1 cm$^2$ són 9 km$^2$ \ $\to$ \ \boxed{27\text{ km}^2}
\end{minipage}
\end{exemple}

\vfill
\begin{minipage}[c]{0.42\linewidth}\centering
\begin{tikzpicture}[scale=0.6,line join=round]
  \newcommand{\cub}[3]{\draw[fmnavy,fill=#3!40] (#1,#2) -- ++(1,0) -- ++(0,1) -- ++(-1,0) -- cycle; \draw[fmnavy,fill=#3!25] (#1,#2+1) -- ++(0.45,0.35) -- ++(1,0) -- ++(-0.45,-0.35) -- cycle; \draw[fmnavy,fill=#3!60] (#1+1,#2) -- ++(0.45,0.35) -- ++(0,1) -- ++(-0.45,-0.35) -- cycle;}
  \cub{0}{0}{fmA}
  \foreach \z in {1,0} \foreach \y in {0,1} \foreach \x in {0,1} {\cub{3.4+\x+0.45*\z}{\y+0.35*\z}{fmB}}
  \node[below,font=\small] at (0.6,-0.2) {$k=1$}; \node[below,font=\small] at (4.6,-0.2) {$k=2$: 8 cubets};
\end{tikzpicture}
\end{minipage}\hfill
\begin{minipage}[c]{0.55\linewidth}
\begin{tcolorbox}[enhanced,colback=fmyellow!18,colframe=fmgold,arc=4mm,boxrule=1.6pt,before upper=\centering,top=2mm,bottom=2mm]
{\Large$\dfrac{L'}{L}=k\qquad\dfrac{A'}{A}=k^2\qquad\dfrac{V'}{V}=k^3$}
\end{tcolorbox}
{\small\raggedright A escala $1:n$: les àrees reals són $n^2$ vegades les del plànol, i una maqueta té un volum $n^3$ vegades menor.}
\end{minipage}

\vfill
\begin{exemple}[Exemple: \ maqueta d'un vaixell a escala 1:100, coberta de 0,6 m$^2$ i casc de 0,04 m$^3$]
$k=100$: \quad $A=0{,}6\cdot100^2=\boxed{6000\text{ m}^2}$ \qquad $V=0{,}04\cdot100^3=\boxed{40\,000\text{ m}^3}$
\end{exemple}

\vfill
\begin{compte}[La mateixa unitat]
Abans de calcular una escala, passa les dues mesures a la \textbf{mateixa unitat}: \ 12 km $=1\,200\,000$ cm.
\end{compte}

\vfill
% ============================================================ resum
\resumtitol
\begin{tcbraster}[raster columns=2,raster equal height=rows,raster column skip=14pt,raster row skip=30pt,raster before skip=16pt]
\rtile{10.1}{Proporcions}{$\frac ab=\frac cd\ \Leftrightarrow\ a\cdot d=b\cdot c$.\\[1ex] Multipliquem en creu.}
\rtile{10.2}{Tales}{Paral·leles sobre dues secants: segments proporcionals.\\[1ex] Recíproc: proporcionals $\Rightarrow$ paral·leles.}
\rtile{10.3}{Dividir un segment}{$n$ marques iguals i paral·leles.\\[1ex] Proporcional: repartiment proporcional.}
\rtile{10.4}{Posició de Tales}{Angle comú i costats paral·lels.\\[1ex] Costats sencers des del vèrtex comú.}
\rtile{10.5}{Semblança}{Angles iguals i costats proporcionals.\\[1ex] Longituds $\times k$, àrees $\times k^2$.}
\rtile{10.6}{Criteris}{Dos angles, o tres costats proporcionals, o dos costats i l'angle comprés.}
\rtile{10.7}{Altures}{Ombres, mirall, estaca: dos triangles semblants i una proporció.}
\rtile{10.8}{Escales}{$1:n$: real $=$ plànol $\times n$.\\[1ex] Volums $\times k^3$.}
\end{tcbraster}
\end{document}
